\clearpage\pagestyle{otMain}\fontsize{11}{15}\selectfont
\corrigesChapitre{2}{Réels et intervalles}
\begin{blocLivre}{cor-2-1}
\begin{corrige}{1}{Le plus petit ensemble}
\demarche{Simplifier le nombre avant de classer son écriture.}
$-4$ est entier relatif mais non naturel : $\mathbb Z$. $7/8=0{,}875$ est décimal non entier : $\mathbb D$. $2/3$ est rationnel non décimal : $\mathbb Q$. $\sqrt{49}=7$ appartient à $\mathbb N$. Enfin $\sqrt3$ est irrationnel : parmi les ensembles proposés, seul $\mathbb R$ le contient.
\end{corrige}
\end{blocLivre}
\begin{blocLivre}{cor-2-2}
\begin{corrige}{2}{Comparer exactement}
\demarche{Utiliser les carrés pour comparer deux nombres positifs.}
$-3/2=-1{,}5<-1{,}4$. Les deux autres nombres sont positifs. Comme $(7/5)^2=49/25=1{,}96<2$, on a $7/5<\sqrt2$. Ainsi $\res{-3/2<-1{,}4<7/5<\sqrt2}$. Comparer les carrés était légitime pour ces deux nombres positifs.
\end{corrige}
\end{blocLivre}
\begin{blocLivre}{cor-2-3}
\begin{corrige}{3}{Du texte à l’intervalle}
\demarche{Faire correspondre chaque inégalité à son crochet.}
Les ensembles sont respectivement $\res{[-2;5[}$, $\res{[3;+\infty[}$ et $\res{]-\infty;-1[}$. Le nombre 5 n’appartient pas au premier : le crochet de droite est ouvert. En revanche, $-2$ lui appartient. Le symbole $+\infty$ n’est pas un nombre réel que l’on pourrait inclure.
\end{corrige}
\end{blocLivre}
\begin{blocLivre}{cor-2-4}
\begin{corrige}{4}{Union et intersection}
\demarche{Repérer la zone commune, puis la zone couverte par au moins un intervalle.}
Les nombres présents dans les deux intervalles vérifient $0<x\le2$ : $\res{I\cap J=]0;2]}$. Les deux intervalles se chevauchent, leur réunion couvre sans trou de $-3$ à 5 : $\res{I\cup J=[-3;5]}$. En particulier 0 reste dans la réunion grâce à $I$, même s’il est exclu de $J$.
\label{fig-cor-2-4}
\figureGeo[corrige]{intervalles}
\end{corrige}
\end{blocLivre}
\begin{blocLivre}{cor-2-5}
\begin{corrige}{5}{Distances sur une droite}
\demarche{Transformer la différence signée en distance positive.}
La distance est la valeur absolue de la différence : $|-4-3|=|-7|=\res7$ ; $|2{,}5-(-1{,}5)|=|4|=\res4$. Échanger les deux points change le signe de la différence mais pas sa valeur absolue. C’est pourquoi une distance reste la même dans les deux sens.
\end{corrige}
\end{blocLivre}
\begin{blocLivre}{cor-2-6}
\begin{corrige}{6}{Un intervalle centré}
\demarche{Lire une inégalité de valeur absolue comme une distance.}
La condition demande les points situés à distance au plus 3 de l’abscisse 2. Les extrémités sont $2-3=-1$ et $2+3=5$. Elles sont autorisées par « au plus ». L’ensemble des solutions est $\res{[-1;5]}$. Le centre vaut bien $(-1+5)/2=2$.
\label{fig-cor-2-6}
\figureGeo[corrige]{intervalle-centre}
\end{corrige}
\end{blocLivre}
\begin{blocLivre}{cor-2-7}
\begin{corrige}{7}{Un encadrement de racine}
\demarche{Contrôler un encadrement par des carrés exacts.}
Les nombres sont positifs. On calcule $2{,}23^2=4{,}9729<5$ et $2{,}24^2=5{,}0176>5$. La fonction carré conserve l’ordre sur les positifs ; l’encadrement est donc établi. Tous les nombres de cet intervalle s’arrondissent à $\res{2{,}2}$ au dixième, car ils sont inférieurs à 2,25.
\end{corrige}
\end{blocLivre}
\begin{blocLivre}{cor-2-8}
\begin{corrige}{8}{Une erreur absolue}
\demarche{Distinguer rayon d’incertitude et amplitude de l’intervalle.}
Notons $L$ la longueur réelle. L’énoncé se traduit par $|L-12{,}4|\le0{,}15$. Donc $\res{L\in[12{,}25;12{,}55]}$ en centimètres. La largeur de l’intervalle est $12{,}55-12{,}25=\res{0{,}30\text{ cm}}$. L’erreur maximale est le rayon, pas la largeur totale.
\end{corrige}
\end{blocLivre}
\begin{blocLivre}{cor-2-9}
\begin{corrige}{9}{Complémentaire d’un intervalle}
\demarche{Examiner les deux côtés et traiter chaque borne séparément.}
Ne pas vérifier $-2<x\le4$ signifie être à gauche, en incluant $-2$, ou strictement à droite de 4. Le complémentaire est donc $\res{]-\infty;-2]\cup]4;+\infty[}$. Le nombre 4 reste exclu du complémentaire puisqu’il appartient à l’intervalle initial.
\end{corrige}
\end{blocLivre}
\begin{blocLivre}{cor-2-10}
\begin{corrige}{10}{Mesures et chiffres significatifs}
\demarche{Calculer avant d’arrondir et adapter la conversion à la grandeur.}
Le produit exact des données est $12{,}3\times4{,}56=56{,}088$ cm$^2$. À trois chiffres significatifs, il devient $\res{56{,}1\text{ cm}^2}$. Comme $1\text{ m}^2=10000\text{ cm}^2$, la valeur non arrondie est $\res{0{,}0056088\text{ m}^2}$. Pour les autres conversions : $\res{750\text{ mL}=0{,}75\text{ L}}$ ; $\res{2{,}5\text{ h}=2\text{ h }30\text{ min}}$ ; $72\times1000/3600=\res{20\text{ m/s}}$ ; $\res{1{,}4\text{ kg}=1400\text{ g}}$ ; $\res{3\text{ dm}^3=3\text{ L}}$. Le facteur dépend de la grandeur : une aire se convertit avec le carré du facteur de longueur.
\end{corrige}
\end{blocLivre}
\begin{blocLivre}{cor-2-11}
\begin{corrige}{11}{Deux contraintes de mesure}
\demarche{Transformer chaque contrainte puis prendre une intersection.}
La première condition donne $8\le x\le12$, soit $[8;12]$. La seconde donne $]9;13]$. Les deux conditions doivent être satisfaites simultanément : on prend leur intersection. La réponse est $\res{x\in]9;12]}$. La borne 9 est exclue, tandis que 12 respecte bien les deux contraintes.
\end{corrige}
\end{blocLivre}
\begin{blocLivre}{cor-2-12}
\begin{corrige}{12}{Un décimal caché}
\demarche{Réduire la fraction avant d’examiner son dénominateur.}
$13/40=325/1000=0{,}325$ est décimal. $21/28=3/4=75/100=0{,}75$ l’est aussi. La fraction $7/12$ est irréductible et son dénominateur contient un facteur 3 ; elle ne peut avoir un dénominateur puissance de 10 après réduction. Elle n’est donc pas décimale. Ce dernier critère est démontré dans le parcours d’approfondissement.
\end{corrige}
\end{blocLivre}
\begin{blocLivre}{cor-2-13}
\begin{corrige}{13}{Un faux calcul sur les racines}
\demarche{Construire un contre-exemple exact et comparer ses carrés.}
Prendre $a=b=4$ donne $\sqrt8=2\sqrt2$ à gauche et $2+2=4$ à droite. Ces nombres sont distincts puisque leurs carrés valent 8 et 16. De plus $\res{\sqrt8<4}$. La racine d’une somme ne se distribue pas ; cette erreur ne doit pas être confondue avec la propriété du produit de deux réels positifs.
\end{corrige}
\end{blocLivre}
\begin{blocLivre}{cor-2-14}
\begin{corrige}{14}{Une amplitude imposée}
\demarche{Propager un encadrement par des opérations qui conservent l’ordre.}
Multiplier par 3 conserve le sens : $5{,}196<3\sqrt3<5{,}199$. Ajouter 2 donne $\res{7{,}196<2+3\sqrt3<7{,}199}$. L’amplitude est $0{,}003$. Tout cet intervalle est entre 7,195 et 7,205 : l’arrondi au centième est donc $\res{7{,}20}$. L’encadrement permet de garantir l’arrondi.
\end{corrige}
\end{blocLivre}
\begin{blocLivre}{cor-2-15}
\begin{corrige}{15}{Balayer pour encadrer}
\demarche{Remplacer des comparaisons décimales par des calculs entiers exacts.}
Le nombre $k/100$ est positif et $(k/100)^2=k^2/10000$. Comparer son carré à 2 revient donc à comparer $k^2$ à 20000. Or $141^2=19881<20000<20164=142^2$. Ainsi $\res{1{,}41<\sqrt2<1{,}42}$. Le recours aux entiers évite ici les erreurs d’arrondi lors du balayage.
\end{corrige}
\end{blocLivre}
